% Fedora WSL printable cheat sheet.
% Source of truth: docs/platforms/fedora-wsl.md, "PATH and Windows
% interoperability", platforms/fedora-wsl/stow/interop/.local/bin/*.
% Regenerate with docs/cheatsheets/generate.sh after editing either.
\documentclass[9pt]{extarticle}
\usepackage{cheatsheet}

\begin{document}
\cssheettitle{Fedora WSL}{Windows is the desktop/terminal host; Fedora WSL is the Linux development runtime}

\begin{multicols}{2}

\cssection{Windows / Linux boundary}
\csnote{The \texttt{[interop]} block in \texttt{/etc/wsl.conf} sets
\texttt{enabled=true} with
\texttt{appendWindowsPath=false}: after \texttt{wsl --shutdown} no Windows
directory leaks into \texttt{PATH}, but an explicit full path to a Windows
executable still runs. There is no KDE/Sway desktop here --- Windows owns the
keyboard.}

\cssection{Windows interop helpers}
\begin{cskeys}{0.34\linewidth}
\csrow{wsl-open <url\textbar path> [...] }{Open existing Linux paths or supported URIs with Windows handlers}
\csrow{wsl-copy}{Send stdin to the Windows clipboard}
\csrow{wsl-paste}{Write the Windows clipboard to stdout}
\end{cskeys}
\csnote{\texttt{BROWSER=wsl-open} is set automatically, so \texttt{gh auth login --web}, \texttt{gx} in LazyVim, and Markdown/LaTeX preview all open Windows's browser/PDF handler. Set \texttt{WINDOWS\_SYSTEM\_ROOT} if Windows is not at \texttt{/mnt/c/Windows}.}

\cssection{Noctty (if used as the Windows terminal)}
% Noctty's own upstream defaults (its docs/getting-started.md): mostly
% Ghostty's non-macOS set, but the split and pane-focus keys are Windows
% specific, so this table is not a copy of the Fedora sheets' Ghostty one.
\begin{cskeys}{0.34\linewidth}
\csrow{Ctrl+Shift+C / V}{Copy / paste}
\csrow{Ctrl+Shift+T / W}{New / close tab}
\csrow{Ctrl+Tab / Ctrl+Shift+Tab}{Next / previous tab}
\csrow{Ctrl+Shift+\textbackslash{} / E}{Split pane right / down}
\csrow{Alt+arrows}{Move between panes}
\csrow{Ctrl+Shift+P}{Command palette}
\csrow{Ctrl+Shift+F}{Start a search}
\csrow{Ctrl+Shift+X}{Copy mode}
\csrow{Ctrl+= / Ctrl+-}{Increase / decrease the font size}
\csrow{Ctrl+Shift+,}{Reload Noctty config (e.g.\ after \texttt{theme})}
\end{cskeys}
\csnote{Noctty shares Ghostty's \texttt{shared.conf}; any other Windows terminal works too. Every row is Noctty's own default --- this repository sets no \texttt{keybind =} --- and \texttt{noctty +list-keybinds} prints the full set. \texttt{Ctrl+Shift+O} splits right as well; Noctty documents the backslash because Narrator can reserve O-based commands.}
\csnote{Keep repositories under \texttt{\textasciitilde/src}, not \texttt{/mnt/c}: metadata, file-watching and I/O are slower and case-insensitive there.}

\renewcommand{\cstermlegend}{\cslegend{Noctty's own default keys are in the Noctty section above; any other Windows terminal uses its own.}}
% Shared terminal/development workflow block, \input by every profile sheet.
% Keep this in sync with docs/reference/keybindings.md; regenerate all sheets together
% (see docs/cheatsheets/README.md) after editing either file.
\cssection{Terminal}
\csnote{No \texttt{keybind =} is set anywhere in this config --- every terminal shortcut is an unmodified upstream default (varies by terminal, platform and version).}
Ghostty on Fedora and macOS; Noctty on Windows for the WSL profile, from the same tracked \texttt{shared.conf}. Either owns tabs/splits/zoom locally; tmux does not duplicate that.\par
% The discovery command and the copy/paste keys differ per platform, so each
% sheet supplies its own line rather than this block naming a binary the WSL
% runtime does not have. cheatsheet.sty defines an empty default, and the
% same goes for \cstermdefaults: the terminal's own default keys differ
% between Ghostty on Linux, Ghostty on macOS and Noctty on Windows, so each
% sheet carries the table that is true of its own terminal.
\cstermlegend
\cstermdefaults

\cssection{Shell (Zsh)}
\begin{cskeys}{0.36\linewidth}
\csrow{ls / ll / la}{eza short / long+git / all}
\csrow{tree}{eza --tree}
\csrow{cat}{bat (syntax highlight)}
\csrow{theme <f>}{switch Catppuccin flavour}
\csrow{tar A.tar.gz P...}{create an archive from P...}
\csrow{untar A.tar.gz}{extract an archive}
\csrow{claude-unsafe}{Claude Code, all permission prompts off (unsafe)}
\csrow{shell-integrations}{report optional tools this shell lacks}
\end{cskeys}
\csnote{<f> is latte, frappe, macchiato or mocha. Terminal, tmux, Starship, bat and Lazygit follow it everywhere; the desktop half needs a platform theme hook.}
\csnote{tar is a shorthand only when the first argument is an archive name; tar -tf, tar -xf, tar --help and command tar behave normally.}

\cssection{Zsh line editing}
\begin{cskeys}{0.25\linewidth}
\csrow{Up / Down}{history entries matching the typed prefix}
\csrow{Ctrl+R}{fuzzy history search (fzf)}
\csrow{Tab}{open/select from the completion menu}
\csrow{Home / End}{beginning / end of line}
\csrow{Delete}{delete the character under the cursor}
\csrow{Ctrl+Left / Right}{move one word backward / forward}
\csrow{\#}{starts an interactive comment}
\end{cskeys}
\csnote{Bindings are taken from terminfo with xterm/application/vt220 fallbacks, so they work in Ghostty, Konsole, WSL, VM consoles, SSH and tmux.}

\cssection{fzf}
\begin{cskeys}{0.20\linewidth}
\csrow{Ctrl+R}{fuzzy history search}
\csrow{Ctrl+T}{fuzzy file selection}
\csrow{Alt+C}{fuzzy cd into a directory}
\end{cskeys}

\cssection{zoxide}
\begin{cskeys}{0.20\linewidth}
\csrow{z name / zi}{jump to a ranked directory / interactive picker}
\end{cskeys}

\cssection{tmux (prefix Ctrl+B, unmodified)}
\begin{cskeys}{0.30\linewidth}
\csrow{tmux new -As X}{create/attach session X}
\csrow{Prefix d}{detach}
\csrow{Prefix s}{choose session}
\csrow{Prefix [}{scroll/copy mode}
\csrow{Prefix ?}{list all key bindings}
\csrow{Mouse scroll / click}{scroll a pane's history, select panes and windows}
\end{cskeys}
\csnote{Thin by design: persistence/remote sessions only; Ghostty stays the local tab/split manager.}

\cssection{Neovim essentials (upstream Vim defaults)}
\begin{cskeys}{0.28\linewidth}
\csrow{i / a / o}{Insert before / after the cursor, open a line below}
\csrow{Esc}{Leave insert mode, back to normal mode}
\csrow{:w / :q / :wq / :q!}{Write / quit / write and quit / quit discarding}
\csrow{h j k l}{Move left / down / up / right}
\csrow{w / b / e}{Next word / previous word / end of word}
\csrow{0 / \$}{Start / end of the line}
\csrow{gg / G}{Top / bottom of the buffer}
\csrow{Ctrl+U / Ctrl+D}{Half a screen up / down}
\csrow{yy / dd / p}{Yank / cut / put a line}
\csrow{ciw / diw}{Change / delete the word under the cursor}
\csrow{v / V / Ctrl+V}{Character / line / block visual mode}
\csrow{/pat, then n / N}{Search, then next / previous match}
\csrow{:\%s/a/b/g}{Replace a with b in the whole buffer}
\csrow{u / Ctrl+R}{Undo / redo}
\csrow{.}{Repeat the last change}
\csrow{gcc / gc + motion}{Toggle a comment on the line / over a motion}
\end{cskeys}
\csnote{Vim's own defaults, not dotfiles bindings: they work in any \texttt{nvim}, including one started with \texttt{--clean} or on a machine this repository has never touched. \texttt{:help index} lists every one of them.}

\cssection{LazyVim / Neovim}
\begin{cskeys}{0.26\linewidth}
\csrow{Space}{WhichKey menu (discover everything)}
\csrow{<lead>ff}{find files}
\csrow{<lead>/}{live grep}
\csrow{<lead>,}{buffers}
\csrow{S-h/S-l}{prev/next buffer}
\csrow{gd / gr}{go to definition / references}
\csrow{K}{hover docs}
\csrow{<lead>ca}{code action}
\csrow{<lead>cr}{rename symbol}
\csrow{<lead>cf}{format buffer}
\csrow{]d / [d}{next/prev diagnostic}
\csrow{<lead>xx}{diagnostics list (Trouble)}
\csrow{<lead>db}{toggle breakpoint (DAP)}
\csrow{<lead>dc}{start/continue debugging}
\csrow{<lead>du}{toggle debug UI}
\csrow{<lead>ft}{floating terminal}
\csrow{<lead>gg}{Lazygit}
\csrow{<lead>e}{file explorer (Snacks), rooted at the project}
\csrow{Ctrl+H/J/K/L}{focus the window left/down/up/right}
\csrow{<lead>- / <lead>\textbar}{split the window below / right}
\csrow{<lead>wd}{close the window}
\csrow{<lead>bd}{close the buffer, keep the window}
\csrow{Ctrl+S}{save the file, from any mode}
\csrow{Alt+J / Alt+K}{move the line (or selection) down / up}
\csrow{s}{Flash: jump to any visible position}
\csrow{<lead>fr}{recent files}
\csrow{<lead>sk}{search every keymap}
\csrow{<lead>l}{Lazy: plugin manager}
\csrow{<lead>qq}{quit all}
\end{cskeys}
\csnote{<lead> is the leader key (Space). These are LazyVim's own defaults, verified against the commit pinned in \texttt{lazy-lock.json}; \texttt{<lead>cf} is the one whose formatter map this repository extends. Full keymap: press Space and read WhichKey, or \texttt{<lead>sk}.}

\cssection{Lazygit}
\begin{cskeys}{0.20\linewidth}
\csrow{lazygit}{open from any Git repo, or <lead>gg in LazyVim}
\csrow{?}{contextual key-binding help for the current panel}
\end{cskeys}

\cssection{Herdr --- agent workspaces (optional AI profile)}
\begin{cskeys}{0.34\linewidth}
\csrow{herdr}{start or reattach the workspace; agents survive a detach}
\csrow{Ctrl+B ?}{list every binding}
\csrow{Ctrl+B c}{new tab}
\csrow{Ctrl+B v / Ctrl+B -}{split the pane right / down}
\csrow{Ctrl+B h/j/k/l}{focus the pane left/down/up/right}
\csrow{Ctrl+B z / Ctrl+B x}{zoom / close the focused pane}
\csrow{Ctrl+B n / Ctrl+B p}{next / previous tab}
\csrow{Ctrl+B 1..9}{jump to tab 1-9}
\csrow{Ctrl+B Shift+N}{new workspace}
\csrow{Ctrl+B w / Ctrl+B g}{workspace picker / goto picker}
\csrow{Ctrl+B b}{toggle the agent sidebar (blocked/working/done/idle)}
\csrow{Ctrl+B [}{copy mode}
\csrow{Ctrl+B q}{detach; the agents keep running}
\end{cskeys}
\csnote{Installed only by \texttt{./install.sh --ai}; these are Herdr's own upstream defaults, and it is mouse-native besides. Its prefix is \texttt{Ctrl+B}, the same key tmux uses --- run one inside the other and the inner multiplexer never sees the prefix. Herdr is for agent panes, tmux for ordinary shell persistence.}


\csfootflow{Full keybindings reference: docs/reference/keybindings.md. Interop config: docs/platforms/fedora-wsl.md. Regenerate: docs/cheatsheets/generate.sh.}

\end{multicols}
\end{document}
